\begin{figure}[htbp]
    \centering
    \shorthandoff{<>}
    \begin{adjustbox}{max width=\linewidth, max height=0.5\textheight, center}
    \begin{tikzpicture}[
        font=\sffamily,
        >=Stealth,
        box/.style={
            draw,
            rectangle,
            fill=blue!8,
            draw=blue!60!black,
            line width=0.8pt,
            minimum width=2cm,
            minimum height=0.95cm,
            align=center,
            rounded corners=3pt,
            font=\scriptsize\bfseries
        },
        lifeline/.style={draw, dashed, color=gray!70, line width=0.6pt},
        msg/.style={font=\tiny, midway, align=center, fill=white, inner sep=1pt},
        note/.style={font=\tiny, align=left, fill=white, inner sep=1pt},
        call/.style={-{Stealth[length=2mm]}, line width=0.7pt, color=black!80},
        resp/.style={{Stealth[length=2mm]}-, dashed, line width=0.7pt, color=black!70}
    ]

        % --- Participantes ---
        \node[box] (M) at (0,0) {App móvil\\\mdseries (M)};
        \node[box] (G) at (2.3,0) {Gateway\\\mdseries (G :8000)};
        \node[box] (A) at (4.8,0) {ASR Service\\\mdseries (A :8001)};
        \node[box] (W) at (7.2,0) {Whisper\\\mdseries (W)};
        \node[box] (P) at (9.3,0) {PLN/NER\\\mdseries (P)};
        \node[box] (S) at (11.4,0) {SNOMED\\\mdseries + CIE-10 (S)};

        % --- Líneas de vida ---
        \pgfmathsetmacro{\lifelineDepth}{-9.8}
        \foreach \n in {M,G,A,W,P,S} {
            \draw[lifeline] (\n.south) -- (\n.center |- 0,\lifelineDepth);
        }
        \node[box, font=\tiny\bfseries] (M bot) at (M.center |- 0,\lifelineDepth) {App móvil};
        \node[box, font=\tiny\bfseries] (G bot) at (G.center |- 0,\lifelineDepth) {Gateway};
        \node[box, font=\tiny\bfseries] (A bot) at (A.center |- 0,\lifelineDepth) {ASR Service};
        \node[box, font=\tiny\bfseries] (W bot) at (W.center |- 0,\lifelineDepth) {Whisper};
        \node[box, font=\tiny\bfseries] (P bot) at (P.center |- 0,\lifelineDepth) {PLN/NER};
        \node[box, font=\tiny\bfseries] (S bot) at (S.center |- 0,\lifelineDepth) {SNOMED};

        % --- Conexión inicial ---
        \draw[call] ($(M.south)+(0,-0.5)$) -- ($(G.south)+(0,-0.5)$)
            node[msg, below] {WSS connect /audio};

        % --- Bucle por fragmento ---
        \coordinate (loopNW) at (-0.5,-0.75);
        \coordinate (loopSE) at (12.1,-9.1);

        \draw[call] ($(M.south)+(0,-1.0)$) -- ($(G.south)+(0,-1.0)$)
            node[msg] {JSON metadata (audio\_chunk)};
        \draw[call] ($(M.south)+(0,-1.55)$) -- ($(G.south)+(0,-1.55)$)
            node[msg] {Binario PCM (audio)};

        \draw[call] ($(G.south)+(0,-2.1)$) -- ($(A.south)+(0,-2.1)$)
            node[msg] {POST /transcribe\\(octet-stream + headers)};
        \draw[call] ($(A.south)+(0,-2.65)$) -- ($(W.south)+(0,-2.65)$)
            node[msg] {Transcripción};
        \draw[resp] ($(W.south)+(0,-3.15)$) -- ($(A.south)+(0,-3.15)$)
            node[msg] {Texto};

        \draw[call] ($(A.south)+(0,-3.65)$) -- ($(P.south)+(0,-3.65)$)
            node[msg] {NER (pln\_medical + pln\_farmacos)};
        \draw[resp] ($(P.south)+(0,-4.15)$) -- ($(A.south)+(0,-4.15)$)
            node[msg] {Entidades};

        \draw[call] ($(A.south)+(0,-4.65)$) -- ($(S.south)+(0,-4.65)$)
            node[msg] {Enriquecimiento SNOMED + CIE-10};
        \draw[resp] ($(S.south)+(0,-5.15)$) -- ($(A.south)+(0,-5.15)$)
            node[msg] {Entidades enriquecidas};

        \draw[call]
            ($(A.south)+(0,-5.55)$) -- ++(0.45,0) -- ++(0,-0.35) -- ($(A.south)+(0,-5.9)$)
            node[note, right, xshift=0.15cm] {Buffer en memoria\\(SessionReportWriter)};

        \draw[resp] ($(G.south)+(0,-6.45)$) -- ($(A.south)+(0,-6.45)$)
            node[msg] {JSON (text, entities, ...)};

        \draw[call]
            ($(G.south)+(0,-6.85)$) -- ++(0.45,0) -- ++(0,-0.35) -- ($(G.south)+(0,-7.2)$)
            node[note, right, xshift=0.15cm] {Acumula \\ sessionEntitiesStore};

        \draw[resp] ($(M.south)+(0,-7.75)$) -- ($(G.south)+(0,-7.75)$)
            node[msg] {Log de transmisión};

        \begin{scope}[on background layer]
            \node[
                draw=orange!80!black,
                fill=orange!8,
                fill opacity=0.25,
                rectangle,
                rounded corners=3pt,
                fit={(loopNW) (loopSE)},
                inner sep=6pt
            ] (loopBox) {};
            \node[
                fill=orange!20,
                draw=orange!80!black,
                font=\tiny\bfseries,
                text=orange!60!black,
                rounded corners=2pt,
                anchor=north west,
                inner sep=2pt
            ] at (loopBox.north west) {LOOP: Por cada fragmento de voz};
        \end{scope}

    \end{tikzpicture}%
    \end{adjustbox}%
    \shorthandon{<>}
    \caption{Diagrama de secuencia del flujo de transcripción en tiempo real en modo \textit{streaming}: captura de audio, transcripción \englishText{ASR} y enriquecimiento semántico por fragmento.}
    \label{fig:diagrama_conceptual_transcripcion}
\end{figure}

\begin{figure}[htbp]
    \centering
    \shorthandoff{<>}
    \begin{adjustbox}{max width=\linewidth, max height=0.55\textheight, center}
    \begin{tikzpicture}[
        font=\sffamily,
        >=Stealth,
        box/.style={
            draw,
            rectangle,
            fill=blue!8,
            draw=blue!60!black,
            line width=0.8pt,
            minimum width=2cm,
            minimum height=0.95cm,
            align=center,
            rounded corners=3pt,
            font=\scriptsize\bfseries
        },
        lifeline/.style={draw, dashed, color=gray!70, line width=0.6pt},
        msg/.style={font=\tiny, midway, align=center, fill=white, inner sep=1pt},
        note/.style={font=\tiny, align=left, fill=white, inner sep=1pt},
        call/.style={-{Stealth[length=2mm]}, line width=0.7pt, color=black!80},
        resp/.style={{Stealth[length=2mm]}-, dashed, line width=0.7pt, color=black!70}
    ]

        % --- Participantes ---
        \node[box] (M) at (0,0) {App móvil\\\mdseries (M)};
        \node[box] (G) at (2.3,0) {Gateway\\\mdseries (G :8000)};
        \node[box] (A) at (4.8,0) {ASR Service\\\mdseries (A :8001)};

        % --- Líneas de vida ---
        \pgfmathsetmacro{\lifelineDepth}{-7.7}
        \foreach \n in {M,G,A} {
            \draw[lifeline] (\n.south) -- (\n.center |- 0,\lifelineDepth);
        }
        \node[box, font=\tiny\bfseries] (M bot) at (M.center |- 0,\lifelineDepth) {App móvil};
        \node[box, font=\tiny\bfseries] (G bot) at (G.center |- 0,\lifelineDepth) {Gateway};
        \node[box, font=\tiny\bfseries] (A bot) at (A.center |- 0,\lifelineDepth) {ASR Service};

        % --- Conexión inicial ---
        \draw[call] ($(M.south)+(0,-0.5)$) -- ($(G.south)+(0,-0.5)$)
            node[msg, below] {WSS connect /audio};

        % --- Bucle por fragmento: solo acumulación, sin llamada al ASR ---
        \coordinate (loopNW) at (-0.5,-0.75);
        \coordinate (loopSE) at (5.9,-2.7);

        \draw[call] ($(M.south)+(0,-1.0)$) -- ($(G.south)+(0,-1.0)$)
            node[msg] {JSON metadata (audio\_chunk)};
        \draw[call] ($(M.south)+(0,-1.55)$) -- ($(G.south)+(0,-1.55)$)
            node[msg] {Binario PCM (audio)};

        \draw[call]
            ($(G.south)+(0,-1.95)$) -- ++(0.45,0) -- ++(0,-0.35) -- ($(G.south)+(0,-2.3)$)
            node[note, right, xshift=0.15cm] {Acumula fragmento en buffer\\(\texttt{batchAudioBuffer.addChunk})};

        \begin{scope}[on background layer]
            \node[
                draw=orange!80!black,
                fill=orange!8,
                fill opacity=0.25,
                rectangle,
                rounded corners=3pt,
                fit={(loopNW) (loopSE)},
                inner sep=6pt
            ] (loopBox) {};
            \node[
                fill=orange!20,
                draw=orange!80!black,
                font=\tiny\bfseries,
                text=orange!60!black,
                rounded corners=2pt,
                anchor=north west,
                inner sep=2pt
            ] at (loopBox.north west) {LOOP: Por cada fragmento de voz (sin enviar al ASR)};
        \end{scope}

        % --- Fin de sesión: vaciado del buffer y única llamada al ASR ---
        \draw[call] ($(M.south)+(0,-3.2)$) -- ($(G.south)+(0,-3.2)$)
            node[msg] {JSON session\_end};

        \draw[call]
            ($(G.south)+(0,-3.6)$) -- ++(0.45,0) -- ++(0,-0.35) -- ($(G.south)+(0,-3.95)$)
            node[note, right, xshift=0.15cm] {\texttt{flushSession}: concatena fragmentos\\(\texttt{Buffer.concat})};

        \draw[call] ($(G.south)+(0,-4.45)$) -- ($(A.south)+(0,-4.45)$)
            node[msg] {POST /transcribe \{batch: true\}\\(audio completo de la sesión)};

        \draw[call]
            ($(A.south)+(0,-4.85)$) -- ++(0.45,0) -- ++(0,-0.35) -- ($(A.south)+(0,-5.2)$)
            node[note, right, xshift=0.15cm] {\textit{Pipeline} interno \\(Mismo que en el caso anterior)};

        \draw[resp] ($(G.south)+(0,-5.7)$) -- ($(A.south)+(0,-5.7)$)
            node[msg] {JSON (text, entities, ...)};

        \draw[call]
            ($(G.south)+(0,-6.1)$) -- ++(0.45,0) -- ++(0,-0.35) -- ($(G.south)+(0,-6.45)$)
            node[note, right, xshift=0.15cm] {Acumula \\ sessionEntitiesStore};

    \end{tikzpicture}%
    \end{adjustbox}%
    \shorthandon{<>}
    \caption{Diagrama de secuencia del flujo de transcripción en modo \textit{batch}: los fragmentos de
        audio se acumulan en memoria durante toda la sesión y se envían como un único bloque al
        servicio \englishText{ASR} al recibir \texttt{session\_end}. A partir de aquí el flujo de
        cierre de sesión es común a ambas modalidades (Figura~\ref{fig:diagrama_conceptual_informe_llm}).}
    \label{fig:diagrama_conceptual_transcripcion_batch}
\end{figure}

\begin{figure}[htbp]
    \centering
    \shorthandoff{<>}
    \begin{adjustbox}{max width=\linewidth, max height=0.55\textheight, center}
    \begin{tikzpicture}[
        font=\sffamily,
        >=Stealth,
        box/.style={
            draw,
            rectangle,
            fill=blue!8,
            draw=blue!60!black,
            line width=0.8pt,
            minimum width=2cm,
            minimum height=0.95cm,
            align=center,
            rounded corners=3pt,
            font=\scriptsize\bfseries
        },
        lifeline/.style={draw, dashed, color=gray!70, line width=0.6pt},
        msg/.style={font=\tiny, midway, align=center, fill=white, inner sep=1.5pt},
        note/.style={font=\tiny, align=left, fill=white, inner sep=1.5pt},
        call/.style={-{Stealth[length=2mm]}, line width=0.7pt, color=black!80},
        resp/.style={{Stealth[length=2mm]}-, dashed, line width=0.7pt, color=black!70}
    ]

        % --- Participantes ---
        \node[box] (M) at (0,0) {App móvil\\\mdseries (M)};
        \node[box] (G) at (2.3,0) {Gateway\\\mdseries (G :8000)};
        \node[box] (A) at (4.8,0) {ASR Service\\\mdseries (A :8001)};
        \node[box] (L) at (7.6,0) {LLM\\\mdseries (Ollama) (L)};

        % --- Líneas de vida ---
        \pgfmathsetmacro{\lifelineDepth}{-8.0}
        \foreach \n in {M,G,A,L} {
            \draw[lifeline] (\n.south) -- (\n.center |- 0,\lifelineDepth);
        }
        \node[box, font=\tiny\bfseries] (M bot) at (M.center |- 0,\lifelineDepth) {App móvil};
        \node[box, font=\tiny\bfseries] (G bot) at (G.center |- 0,\lifelineDepth) {Gateway};
        \node[box, font=\tiny\bfseries] (A bot) at (A.center |- 0,\lifelineDepth) {ASR Service};
        \node[box, font=\tiny\bfseries] (L bot) at (L.center |- 0,\lifelineDepth) {LLM};

        % --- Cierre de sesión ---
        \draw[call] ($(M.south)+(0,-0.5)$) -- ($(G.south)+(0,-0.5)$)
            node[msg, below=4pt] {JSON session\_end};

        % Self-call: espera transcripciones (loop ensanchado, nota debajo para no chocar con el opt)
        \draw[call]
            ($(G.south)+(0,-1.0)$) -- ++(0.6,0) -- ++(0,-0.4) -- ($(G.south)+(0,-1.4)$);
        \node[note, anchor=west] at ($(G.south)+(0.7,-1.2)$) {Espera transcripciones\\pendientes};

        \draw[call] ($(G.south)+(0,-2.15)$) -- ($(A.south)+(0,-2.15)$)
            node[msg, below=4pt] {POST /finalize-session/[id]};

        % --- opt frame: dibujado ANTES del self-call interno para que quepa con holgura ---
        \begin{scope}[on background layer]
            \coordinate (optReportsNW) at (3.5,-2.3);
            \coordinate (optReportsSE) at (6.3,-3.6);
            \node[
                draw=blue!70!black,
                fill=blue!8,
                fill opacity=0.25,
                rectangle,
                rounded corners=3pt,
                fit={(optReportsNW) (optReportsSE)},
                inner sep=4pt
            ] (optReportsBox) {};
            \node[
                fill=blue!20,
                draw=blue!70!black,
                font=\tiny\bfseries,
                text=blue!60!black,
                rounded corners=2pt,
                anchor=north west,
                inner sep=2pt
            ] at ($(optReportsBox.north west)+(0,0.02)$) {opt [\texttt{reports.enabled}]};
        \end{scope}

        % Self-call: escribe Reporte.txt (dentro del opt, con espacio respecto al label del opt)
        \draw[call]
            ($(A.south)+(0,-2.75)$) -- ++(0.6,0) -- ++(0,-0.4) -- ($(A.south)+(0,-3.15)$);
        \node[note, anchor=west] at ($(A.south)+(0.7,-2.95)$) {Escribe \texttt{Reporte.txt}\\en disco};

        \draw[resp] ($(M.south)+(0,-4.0)$) -- ($(G.south)+(0,-4.0)$)
            node[msg, below=4pt] {WS session\_end};

        % --- Informe LLM ---
        \draw[call] ($(G.south)+(0,-4.7)$) -- ($(A.south)+(0,-4.7)$)
            node[msg, below=4pt] {POST /generate-report/[id]};
        \draw[call] ($(A.south)+(0,-5.4)$) -- ($(L.south)+(0,-5.4)$)
            node[msg] {De los datos acumulados, genera informe};
        \draw[resp] ($(L.south)+(0,-6.1)$) -- ($(A.south)+(0,-6.1)$)
            node[msg, below=4pt] {Informe.txt};

        % \draw[resp] ($(G.south)+(0,-6.8)$) -- ($(A.south)+(0,-6.8)$)
        %     node[msg] {JSON report};
        % \draw[resp] ($(M.south)+(0,-7.5)$) -- ($(G.south)+(0,-7.5)$)
        %     node[msg] {WS llm\_report};

    \end{tikzpicture}%
    \end{adjustbox}%
    \shorthandon{<>}
    \caption{Diagrama de secuencia del cierre de sesión y generación del informe clínico mediante \englishText{LLM}.}
    \label{fig:diagrama_conceptual_informe_llm}
\end{figure}

\begin{figure}[htbp]
    \centering
    \shorthandoff{<>}
    \begin{adjustbox}{max width=\linewidth, max height=0.5\textheight, center}
    \begin{tikzpicture}[
        font=\sffamily,
        >=Stealth,
        box/.style={
            draw,
            rectangle,
            fill=blue!8,
            draw=blue!60!black,
            line width=0.8pt,
            minimum width=2cm,
            minimum height=0.95cm,
            align=center,
            rounded corners=3pt,
            font=\scriptsize\bfseries
        },
        lifeline/.style={draw, dashed, color=gray!70, line width=0.6pt},
        msg/.style={font=\tiny, midway, align=center, fill=white, inner sep=1pt},
        note/.style={font=\tiny, align=left, fill=white, inner sep=1pt},
        call/.style={-{Stealth[length=2mm]}, line width=0.7pt, color=black!80},
        resp/.style={{Stealth[length=2mm]}-, dashed, line width=0.7pt, color=black!70}
    ]

        % --- Participantes ---
        \node[box] (M) at (0,0) {App móvil\\\mdseries (M)};
        \node[box] (G) at (2.3,0) {Gateway\\\mdseries (G :8000)};
        \node[box] (A) at (4.8,0) {ASR Service\\\mdseries (A :8001)};
        \node[box] (F) at (7.2,0) {FHIR Gen\\\mdseries (F)};

        % --- Líneas de vida ---
        \pgfmathsetmacro{\lifelineDepth}{-4.5}
        \foreach \n in {M,G,A,F} {
            \draw[lifeline] (\n.south) -- (\n.center |- 0,\lifelineDepth);
        }
        \node[box, font=\tiny\bfseries] (M bot) at (M.center |- 0,\lifelineDepth) {App móvil};
        \node[box, font=\tiny\bfseries] (G bot) at (G.center |- 0,\lifelineDepth) {Gateway};
        \node[box, font=\tiny\bfseries] (A bot) at (A.center |- 0,\lifelineDepth) {ASR Service};
        \node[box, font=\tiny\bfseries] (F bot) at (F.center |- 0,\lifelineDepth) {FHIR Gen};

        % --- Informe FHIR ---
        \draw[call] ($(G.south)+(0,-0.4)$) -- ($(A.south)+(0,-0.4)$)
            node[msg, below=4pt] {POST /generate-fhir-report/[id] + entities};
        \draw[call] ($(A.south)+(0,-0.9)$) -- ($(F.south)+(0,-0.9)$)
            node[msg, below=4pt] {Construye Bundle FHIR};
        \draw[resp] ($(F.south)+(0,-1.4)$) -- ($(A.south)+(0,-1.4)$)
            node[msg, below=4pt] {Fhir\_Reporte.json};

        \draw[resp] ($(G.south)+(0,-1.9)$) -- ($(A.south)+(0,-1.9)$)
            node[msg] {JSON fhir\_report};
        \draw[resp] ($(M.south)+(0,-2.4)$) -- ($(G.south)+(0,-2.4)$)
            node[msg] {WS fhir\_report};

        \draw[call]
            ($(G.south)+(0,-2.8)$) -- ++(0.45,0) -- ++(0,-0.35) -- ($(G.south)+(0,-3.15)$)
            node[note, right, xshift=0.15cm] {Limpia entities\\en memoria};

    \end{tikzpicture}%
    \end{adjustbox}%
    \shorthandon{<>}
    \caption{Diagrama de secuencia de la generación del informe estructurado \englishText{FHIR}.}
    \label{fig:diagrama_conceptual_fhir}
\end{figure}
